% ============================================
% LUC-03 (EN)
% English translation — SPEC-935-R355
% ============================================
\section*{\textbf{LUC-03} --- The Ruins}
\markboth{LUC-03 --- The Ruins}{}

\noindent \textit{The Colony abandons}

\noindent \textbf{March 2026 --- Barbacena, Minas Gerais}

\vspace{0.5em}\noindent\rule{\textwidth}{0.5pt}\vspace{0.5em}

\noindent The Hospital-Colony of Barbacena was closed in 1980, one year after Joaquim's death.

\noindent Today, what remains is a group of abandoned buildings, surrounded by tall weeds and silence. The walls still stand, but the roof has collapsed in several wards. The wooden floor has rotted. The rusted iron beds are still there, some with remains of mattresses, others with nothing.

\noindent Lúcia arrived on a Saturday morning. The sun was strong, but the weeds were too tall for the sun to reach the ground. She cleared her way with her hands, pushing aside dry branches and spider webs.
\noindent Today, what remains is a set of abandoned buildings, surrounded by tall brush and silence. The walls are still standing, but the roof has collapsed in several wards. The wooden floors have rotted. The rusted iron beds are still there, some with remnants of mattresses, others with nothing.

\noindent The ward of the silenced ones was in the back.

\noindent She recognized it from the description in the record: \textit{\textquotedblleft{}End of the hallway. Iron door. Twelve cells.\textquotedblright{}}

\noindent The iron door was ajar. Lúcia pushed it. The creak echoed through the empty hallway.

\noindent Inside, the smell was of wet earth and mold. Beneath it, almost imperceptible, there was a sweetness of burst fruit at the bottom of a can.

\noindent She turned on her phone's flashlight.

\noindent The cells were small. Two meters by three. In each one, a concrete bed. In some, remains of blankets --- gray rags that had fallen apart with time.

\noindent Cell 1,260 was at the end of the hallway.

\noindent Lúcia stopped in front of it.

\noindent The door was of solid wood, with an opening the size of a plate --- the hole through which the food was passed. She peered in.

\noindent Inside, the cell was empty.
Almost.

\noindent On the floor, in the right corner, there was a notebook.

\noindent It was not the diary. The diary was in Lúcia's backpack. It was another notebook. Small, with a black cover and yellowed pages.
\noindent She opened the door. She went in.

\noindent The right eye was lighter. Yellowish.

\noindent It came through her tongue.
\noindent It did not come through her nose.

\noindent She bent down and picked up the notebook. On the cover, written by hand:
\noindent The rotten sweetness thickened there, sticking to the back of her tongue.

\noindent \textit{\textquotedblleft{}Joaquim --- 1945 to 1956\textquotedblright{}}

\noindent The years of silence.

\noindent She opened it. The handwriting was tiny, cramped, as if he wanted to fit as many words as possible on each page.

\noindent The first line said:

\noindent \textit{\textquotedblleft{}She will not let me speak. When I open my mouth, she squeezes my throat. The creature does not want me to tell. But she cannot stop me from writing. Writing is not speaking. Writing is thinking on paper. She does not read my thoughts. She only reads what I write. But she is learning.\textquotedblright{}}

\noindent The last line said:

\noindent \textit{\textquotedblleft{}1956. Today she let me speak. I asked why. She said: 'Because you have already written everything. Now the word is on the paper. It no longer needs to come out through the mouth.' I will speak again. But first I will write more. Writing is the only thing that is still mine.\textquotedblright{}}

\noindent Lúcia closed the notebook.

\noindent The sweetish smell was stronger.

\noindent She raised her eyes.

\noindent In the corner of the cell, on the wall, someone had written something. With ink. Or with something else.

\noindent \textquotedblleft{\textit{Welcome, Dr. Lúcia. You took your time.\textquotedblright{}}}

\noindent The handwriting was the same as the notebook's.

\noindent Joaquim's.

\noindent Written in 1956.

\noindent Foreseeing that she would be there. 70 years later.

\noindent Lúcia took a deep breath.

\noindent The sweetish smell filled her lungs.

\vspace{0.5em}
\noindent\textit{\textbf{↪ Routes:} \rota{LUC-04} \rotasep \rota{MEM-08}}
